% =============================================================================
% paper54_symplectic_geometry_en.tex
% Symplectic Geometry: From Classical Mechanics to Floer Homology
% Author: Michael Fuhrmann
% Build: 139
% Date: 2026-03-12
% =============================================================================

\documentclass[12pt,a4paper]{amsart}
\usepackage{amsmath,amssymb,amsthm}
\usepackage{mathtools}
\usepackage[utf8]{inputenc}
\usepackage[T1]{fontenc}
\usepackage{hyperref}
\usepackage{enumitem,booktabs,array}
\usepackage{geometry}
\geometry{margin=2.5cm}

\newtheorem{theorem}{Theorem}[section]
\newtheorem{lemma}[theorem]{Lemma}
\newtheorem{corollary}[theorem]{Corollary}
\newtheorem{proposition}[theorem]{Proposition}
\newtheorem{conjecture}[theorem]{Conjecture}
\theoremstyle{definition}
\newtheorem{definition}[theorem]{Definition}
\newtheorem{example}[theorem]{Example}
\newtheorem{remark}[theorem]{Remark}

\begin{document}
\title{Symplectic Geometry:\\
From Classical Mechanics to Floer Homology}

\author{Michael Fuhrmann}
\address{Specialist Mathematics Project, Build 139}
\email{specialist-maths@localhost}
\date{\today}

% =============================================================================
% ABSTRACT  (must appear BEFORE \maketitle in amsart)
% =============================================================================
\begin{abstract}
Symplectic geometry is the mathematical framework underlying classical
Hamiltonian mechanics and has grown into one of the most active areas of
modern mathematics, with deep connections to topology, algebra, and
theoretical physics.  Starting from the standard symplectic vector space
$(\mathbb{R}^{2n}, \omega_{\mathrm{std}})$, we develop the theory of
symplectic manifolds, state and prove Darboux's theorem (every symplectic
manifold is locally standard), and derive the Hamiltonian equations of
motion.  We then treat Liouville's theorem on phase-space volume preservation
and Poincar\'{e}'s recurrence theorem---both proved results.  Lagrangian
submanifolds, the Arnold conjecture (open in general; proved for tori),
moment maps, and the Marsden--Weinstein symplectic reduction theorem are
discussed in detail.  The paper culminates with Gromov's celebrated
non-squeezing theorem (1985) and an introduction to Floer homology.
Applications range from integrable systems and KAM theory to geometric
quantisation.
\end{abstract}

\maketitle

\tableofcontents

% =============================================================================
\section{Introduction}
% =============================================================================

Symplectic geometry originated in the nineteenth century as the geometric
language of classical mechanics.  The word \emph{symplectic} was coined by
Hermann Weyl in his 1939 monograph \emph{The Classical Groups} as a Greek
translation of the Latin \emph{complex}; it describes the intertwined
structure of positions and momenta in phase space.

At its heart lies a single mathematical object: a closed, non-degenerate
differential 2-form $\omega$ on a smooth even-dimensional manifold $M$.
This deceptively simple axiom encodes an enormous amount of geometry and
dynamics.  Unlike Riemannian geometry, which features a rich local structure
(curvature), Darboux's theorem shows that symplectic geometry is \emph{locally
flat}: every symplectic manifold looks the same near every point.  All
interesting phenomena are therefore of a \emph{global} nature, and this is
precisely what makes the field both challenging and beautiful.

\subsection*{Historical overview}

The roots of symplectic geometry lie in the work of Lagrange and Hamilton on
analytical mechanics (1788--1835).  Jacobi recognised the invariance of the
Poisson bracket under canonical transformations.  Poincar\'{e}'s recurrence
theorem (1890) was the first purely topological result in the subject.  The
modern era begins with Eliashberg and Gromov's discovery of \emph{symplectic
rigidity}: despite the absence of local invariants, there exist surprising
global obstructions to symplectic embeddings---most strikingly captured by
Gromov's non-squeezing theorem (1985).

\subsection*{Organisation of the paper}

Section~\ref{sec:manifolds} introduces symplectic manifolds and their
canonical examples.  Darboux's theorem is stated and proved in
Section~\ref{sec:darboux}.  Hamiltonian mechanics occupies
Section~\ref{sec:hamilton}.  Symplectomorphisms, Liouville, and Poincar\'{e}
are treated in Section~\ref{sec:symplecto}.  Lagrangian submanifolds and the
Arnold conjecture are covered in Section~\ref{sec:lagrangian}.  Moment maps
and reduction appear in Section~\ref{sec:moment}.  Gromov's non-squeezing
theorem is Section~\ref{sec:gromov}.  Floer homology is introduced in
Section~\ref{sec:floer}.  Applications conclude the paper in
Section~\ref{sec:applications}.

% =============================================================================
\section{Symplectic Manifolds}
\label{sec:manifolds}
% =============================================================================

\begin{definition}
A \emph{symplectic manifold} is a pair $(M, \omega)$ where $M$ is a smooth
manifold (without boundary) and $\omega \in \Omega^{2}(M)$ is a differential
2-form satisfying
\begin{enumerate}[label=(\roman*)]
  \item \textbf{Non-degeneracy:}  $\omega_p(u,v) = 0$ for all $v \in T_p M$
        implies $u = 0$.
  \item \textbf{Closedness:}  $d\omega = 0$.
\end{enumerate}
\end{definition}

Non-degeneracy forces $M$ to be even-dimensional, say $\dim M = 2n$.  The
$n$-th exterior power $\omega^n := \omega \wedge \cdots \wedge \omega$
($n$ factors) is then a nowhere-vanishing volume form, so every symplectic
manifold is orientable.

\subsection{The standard symplectic form on \texorpdfstring{$\mathbb{R}^{2n}$}{R2n}}

Choose coordinates $(q_1, \ldots, q_n, p_1, \ldots, p_n)$ on $\mathbb{R}^{2n}$.
The \emph{standard symplectic form} is
\begin{equation}\label{eq:omega_std}
  \omega_{\mathrm{std}} = \sum_{i=1}^{n} dq_i \wedge dp_i.
\end{equation}
In matrix form, for vectors $u, v \in \mathbb{R}^{2n}$:
\begin{equation}
  \omega_{\mathrm{std}}(u, v) = u^{\top} J\, v, \qquad
  J = \begin{pmatrix} 0 & I_n \\ -I_n & 0 \end{pmatrix}.
\end{equation}
One verifies $d\omega_{\mathrm{std}} = 0$ (constant coefficients) and
$\det J = 1 \neq 0$ (non-degenerate).

\subsection{Cotangent bundles}

Let $Q$ be any smooth $n$-manifold (configuration space).  The cotangent
bundle $T^*Q$ carries a canonical 1-form $\lambda$ (\emph{Liouville form},
also called \emph{tautological form}): if $\pi: T^*Q \to Q$ is the projection
and $\alpha \in T^*Q$, then
\[
  \lambda_\alpha(v) := \alpha\bigl(d\pi_\alpha(v)\bigr), \quad v \in T_\alpha(T^*Q).
\]
The form $\omega := -d\lambda$ is the canonical symplectic form on $T^*Q$.
In local coordinates $(q^i, p_i)$:
\[
  \lambda = \sum_i p_i \, dq^i, \qquad \omega = \sum_i dq^i \wedge dp_i.
\]
Cotangent bundles are the prototypical symplectic manifolds; they provide the
phase space for any mechanical system with configuration space $Q$.

\subsection{Further examples}

\begin{example}[Surfaces]
Every oriented surface $\Sigma$ with an area form $\sigma$ (nowhere-vanishing
2-form) is a symplectic manifold $(\Sigma, \sigma)$, since $d\sigma = 0$
automatically for dimension reasons.
\end{example}

\begin{example}[K\"{a}hler manifolds]
A K\"{a}hler manifold $(M, g, J, \omega)$ is simultaneously Riemannian,
complex, and symplectic.  The symplectic form $\omega(u,v) = g(Ju, v)$ is
closed by the K\"{a}hler identity.  Complex projective space $\mathbb{CP}^n$
with the Fubini--Study form is the primary example.
\end{example}

\begin{example}[Coadjoint orbits]
Let $G$ be a Lie group with Lie algebra $\mathfrak{g}$ and let
$\mu \in \mathfrak{g}^*$.  The coadjoint orbit
$\mathcal{O}_\mu = \{\mathrm{Ad}^*_g \mu : g \in G\}$ carries the
Kirillov--Kostant--Souriau (KKS) symplectic form:
\[
  \omega_\mu(\xi_M, \eta_M) = \mu([\xi, \eta]), \quad \xi, \eta \in \mathfrak{g},
\]
where $\xi_M$ denotes the fundamental vector field of $\xi$.
\end{example}

% =============================================================================
\section{Darboux's Theorem}
\label{sec:darboux}
% =============================================================================

One of the most fundamental results in symplectic geometry distinguishes it
sharply from Riemannian geometry: there are no local invariants.

\begin{theorem}[Darboux, 1882]\label{thm:darboux}
Let $(M, \omega)$ be a $2n$-dimensional symplectic manifold and let
$p \in M$.  Then there exists a neighbourhood $U \ni p$ and local coordinates
$(q_1, \ldots, q_n, p_1, \ldots, p_n)$ on $U$ such that
\[
  \omega\big|_U = \sum_{i=1}^{n} dq_i \wedge dp_i.
\]
Such coordinates are called \emph{Darboux coordinates} or \emph{canonical
coordinates}.
\end{theorem}

\begin{proof}
We follow the Moser path method.  Choose any coordinate chart around $p$;
after a linear change of coordinates we may assume we are working on an open
subset $U \subset \mathbb{R}^{2n}$ with $p = 0$, and that
$\omega_0 := \omega_p = \omega_{\mathrm{std},0}$.

Define a smooth family of 2-forms
\[
  \omega_t = (1-t)\,\omega_{\mathrm{std}} + t\,\omega, \quad t \in [0,1].
\]
Since $\omega_{\mathrm{std}}$ and $\omega$ agree at the origin, $\omega_t$ is
non-degenerate on a (possibly shrunk) neighbourhood $U$ for all $t \in [0,1]$.
Each $\omega_t$ is also closed:
\[
  d\omega_t = (1-t)\,d\omega_{\mathrm{std}} + t\,d\omega = 0.
\]

We seek a time-dependent vector field $X_t$ on $U$ whose flow $\varphi_t$
satisfies $\varphi_t^* \omega_t = \omega_{\mathrm{std}}$.  Differentiating:
\[
  0 = \frac{d}{dt}(\varphi_t^* \omega_t) = \varphi_t^*\!\left(
    \mathcal{L}_{X_t}\omega_t + \dot{\omega}_t \right),
\]
so we need $\mathcal{L}_{X_t}\omega_t + \dot{\omega}_t = 0$, i.e.,
\begin{equation}\label{eq:moser}
  \iota_{X_t} \omega_t + \tau = 0, \quad \tau := \dot{\omega}_t = \omega - \omega_{\mathrm{std}}.
\end{equation}
The form $\tau = \omega - \omega_{\mathrm{std}}$ is closed; by the Poincar\'{e}
lemma on a star-shaped neighbourhood, $\tau = d\sigma$ for some 1-form
$\sigma$ that can be chosen to vanish at $0$ (since $\tau_0 = 0$).
Equation~\eqref{eq:moser} then becomes $\iota_{X_t}\omega_t = -\sigma$,
which has a unique solution $X_t$ because $\omega_t$ is non-degenerate.
The flow $\varphi_1$ of $X_t$ is the desired Darboux coordinate map.
\end{proof}

\begin{remark}
The Darboux theorem has no analogue in Riemannian geometry: not every
Riemannian manifold is locally flat (curvature obstructs this).  In
symplectic geometry, closedness of $\omega$ plays the role that ``curvature
zero'' plays in the Riemannian setting.
\end{remark}

\begin{corollary}
The only local invariant of a symplectic manifold is its dimension.
\end{corollary}

% =============================================================================
\section{Hamiltonian Mechanics}
\label{sec:hamilton}
% =============================================================================

\subsection{The Hamiltonian vector field}

Let $(M, \omega)$ be a symplectic manifold.  Given a smooth function
$H: M \to \mathbb{R}$ (the \emph{Hamiltonian}), the non-degeneracy of $\omega$
guarantees the existence of a unique vector field $X_H$ satisfying
\begin{equation}\label{eq:ham_vf}
  \iota_{X_H}\omega = dH, \quad \text{i.e.,} \quad
  \omega(X_H, \cdot) = dH(\cdot).
\end{equation}
The vector field $X_H$ is called the \emph{Hamiltonian vector field} of $H$.

\subsection{Hamilton's equations}

In Darboux coordinates $(q_i, p_i)$, equation~\eqref{eq:ham_vf} gives
\begin{equation}\label{eq:hamilton}
  \dot{q}_i = \frac{\partial H}{\partial p_i}, \qquad
  \dot{p}_i = -\frac{\partial H}{\partial q_i}, \quad i = 1, \ldots, n.
\end{equation}
These are \emph{Hamilton's equations of motion}.  They describe the
evolution of a mechanical system with $n$ degrees of freedom.

\subsection{Energy conservation}

\begin{proposition}
The Hamiltonian $H$ is preserved along the flow of $X_H$.
\end{proposition}
\begin{proof}
$\dot{H} = dH(X_H) = \omega(X_H, X_H) = 0$, by antisymmetry of $\omega$.
\end{proof}

\subsection{The Poisson bracket}

The \emph{Poisson bracket} of two smooth functions $f, g \in C^\infty(M)$ is
\begin{equation}\label{eq:poisson}
  \{f, g\} := \omega(X_f, X_g) = \sum_{i=1}^{n}
  \left(\frac{\partial f}{\partial q_i}\frac{\partial g}{\partial p_i}
  - \frac{\partial f}{\partial p_i}\frac{\partial g}{\partial q_i}\right).
\end{equation}

\begin{proposition}
$(C^\infty(M), \{\cdot,\cdot\})$ is a Lie algebra.  In particular:
\begin{enumerate}[label=(\roman*)]
  \item \textbf{Antisymmetry:} $\{f,g\} = -\{g,f\}$.
  \item \textbf{Jacobi identity:} $\{f,\{g,h\}\} + \{g,\{h,f\}\} + \{h,\{f,g\}\} = 0$.
  \item \textbf{Leibniz rule:} $\{fg, h\} = f\{g,h\} + g\{f,h\}$.
\end{enumerate}
\end{proposition}

\subsection{Example: the harmonic oscillator}

The harmonic oscillator with $n = 1$ has Hamiltonian
\[
  H(q, p) = \frac{p^2 + \omega^2 q^2}{2}.
\]
Hamilton's equations give $\dot{q} = p$ and $\dot{p} = -\omega^2 q$, yielding
circular orbits in phase space $(q, p)$ of radius
$r = \sqrt{2H/\omega^2}$.  The period is $T = 2\pi/\omega$ independent of
amplitude---a direct consequence of the linearity of the equations.

% =============================================================================
\section{Symplectomorphisms}
\label{sec:symplecto}
% =============================================================================

\begin{definition}
A \emph{symplectomorphism} between symplectic manifolds $(M_1, \omega_1)$ and
$(M_2, \omega_2)$ is a diffeomorphism $\varphi: M_1 \to M_2$ satisfying
$\varphi^* \omega_2 = \omega_1$.
\end{definition}

The flow $\varphi_t$ of any Hamiltonian vector field $X_H$ consists of
symplectomorphisms.  This follows from
$\frac{d}{dt}\varphi_t^*\omega = \varphi_t^*(\mathcal{L}_{X_H}\omega) = 0$,
since $\mathcal{L}_{X_H}\omega = \iota_{X_H}d\omega + d(\iota_{X_H}\omega)
= 0 + d(dH) = 0$.

\subsection{Liouville's theorem}

\begin{theorem}[Liouville, 1838]\label{thm:liouville}
Every Hamiltonian flow preserves the symplectic volume form $\omega^n$.
Equivalently, the phase-space volume of any measurable set is preserved under
the Hamiltonian flow.
\end{theorem}

\begin{proof}
Since $\varphi_t^*\omega = \omega$, we have
\[
  \varphi_t^*(\omega^n) = (\varphi_t^*\omega)^n = \omega^n.
\]
The volume of a Borel set $A \subset M$ under the flow is
$\mathrm{Vol}(\varphi_t(A)) = \int_{\varphi_t(A)} \omega^n
= \int_A \varphi_t^*\omega^n = \int_A \omega^n = \mathrm{Vol}(A)$.
\end{proof}

\begin{remark}
This is the celebrated Liouville theorem of statistical mechanics.  It
implies that the density of states in phase space is conserved---the
foundation of the microcanonical ensemble in statistical mechanics.
\end{remark}

\subsection{Poincar\'{e}'s recurrence theorem}

\begin{theorem}[Poincar\'{e}, 1890]\label{thm:poincare}
Let $\varphi: M \to M$ be a volume-preserving map on a finite-measure
space $(M, \mu)$.  Then for every measurable set $A$ with $\mu(A) > 0$ and
for $\mu$-almost every $x \in A$, there exists $N \geq 1$ such that
$\varphi^N(x) \in A$.
\end{theorem}

\begin{proof}
Define $B := \bigcup_{k=1}^{\infty} \varphi^{-k}(A^c \cap A)$.  One shows
that $B$ has measure zero by the following counting argument: the sets
$\varphi^{-k}(A)$ for $k \geq 0$ must overlap for almost all points, since
otherwise their total measure would exceed $\mu(M) < \infty$.  Formally:
the sets $A \setminus \bigcup_{k=1}^N \varphi^{-k}(A)$ have measure
decreasing to zero as $N \to \infty$ (by volume preservation and finiteness
of total measure), so the set of points in $A$ that never return has
measure zero.
\end{proof}

% =============================================================================
\section{Lagrangian Submanifolds}
\label{sec:lagrangian}
% =============================================================================

\begin{definition}
A submanifold $L \subset (M^{2n}, \omega)$ is called
\begin{enumerate}[label=(\roman*)]
  \item \emph{isotropic} if $\omega|_L = 0$ and $\dim L \leq n$,
  \item \emph{Lagrangian} if $\omega|_L = 0$ and $\dim L = n$.
\end{enumerate}
\end{definition}

Lagrangian submanifolds are the maximally isotropic submanifolds; they have
the largest possible dimension compatible with the isotropy condition.

\begin{example}[Zero section]
In $T^*Q$, the zero section $Q \subset T^*Q$ (where all momenta vanish) is
a Lagrangian submanifold.  Indeed, $\omega|_Q = d\lambda|_Q = 0$ since
$\lambda = p\,dq$ vanishes on the zero section.
\end{example}

\begin{example}[Graph of a closed 1-form]
If $\alpha \in \Omega^1(Q)$ is closed, $d\alpha = 0$, then its graph
$\Gamma_\alpha = \{(q, \alpha_q) : q \in Q\} \subset T^*Q$
is a Lagrangian submanifold.  In particular, graphs of exact forms $\alpha = df$
(for $f \in C^\infty(Q)$) are always Lagrangian.
\end{example}

\subsection{The Arnold conjecture}

A major open problem in symplectic topology concerns the minimum number of
fixed points of Hamiltonian diffeomorphisms.

\begin{conjecture}[Arnold, 1965]\label{conj:arnold}
Let $(M, \omega)$ be a closed symplectic manifold and let $\varphi \in
\mathrm{Ham}(M, \omega)$ be a Hamiltonian diffeomorphism.  Then $\varphi$
has at least as many fixed points as a smooth function on $M$ has critical
points, i.e.,
\[
  \#\mathrm{Fix}(\varphi) \geq \sum_{k=0}^{2n} \beta_k(M),
\]
where $\beta_k = \dim H_k(M; \mathbb{Z}/2)$ are the Betti numbers of $M$.
\end{conjecture}

\begin{remark}
This conjecture is \textbf{open in general}.  It has been proved in special
cases:
\begin{itemize}
  \item \textit{Tori:} Conley and Zehnder proved the Arnold conjecture for
    $M = \mathbb{T}^{2n}$ in 1983.  A Hamiltonian diffeomorphism of the
    $2n$-torus has at least $2n+1$ fixed points (and at least $2^{2n}$
    non-degenerate ones).
  \item \textit{Surfaces:} Proved for all closed surfaces.
  \item \textit{Floer homology approach:} Floer (1989) proved the conjecture
    for monotone manifolds using his homology theory.
\end{itemize}
The general case for arbitrary closed symplectic manifolds remains open.
\end{remark}

% =============================================================================
\section{Moment Maps}
\label{sec:moment}
% =============================================================================

\subsection{Group actions and moment maps}

Let $G$ be a Lie group acting on $(M, \omega)$ by symplectomorphisms
(a \emph{symplectic action}).  For $\xi \in \mathfrak{g}$, let $\xi_M$
denote the fundamental vector field: $\xi_M(x) = \frac{d}{dt}\big|_{t=0}
\exp(t\xi) \cdot x$.

\begin{definition}
The symplectic action of $G$ on $(M, \omega)$ is called \emph{Hamiltonian}
if there exists a map $\mu: M \to \mathfrak{g}^*$ (the \emph{moment map})
satisfying
\begin{enumerate}[label=(\roman*)]
  \item $\iota_{\xi_M}\omega = d\langle \mu, \xi \rangle$ for all
        $\xi \in \mathfrak{g}$ (i.e., $\xi_M = X_{\langle\mu,\xi\rangle}$),
  \item $\mu$ is equivariant: $\mu(g \cdot x) = \mathrm{Ad}^*_g \mu(x)$.
\end{enumerate}
\end{definition}

\begin{example}[$G = S^1$ acting on $\mathbb{C}^n$]
The $S^1$-action $e^{i\theta}(z_1, \ldots, z_n) = (e^{i\theta}z_1, \ldots,
e^{i\theta}z_n)$ on $(\mathbb{C}^n, \frac{i}{2}\sum dz_k \wedge d\bar{z}_k)$
has moment map $\mu(z) = -\frac{1}{2}|z|^2$.
\end{example}

\subsection{Marsden--Weinstein reduction}

\begin{theorem}[Marsden--Weinstein, 1974]\label{thm:mw}
Let $G$ act on $(M, \omega)$ in a Hamiltonian fashion with moment map
$\mu: M \to \mathfrak{g}^*$.  Let $\lambda \in \mathfrak{g}^*$ be a regular
value of $\mu$ and assume $G_\lambda := \{\,g \in G : \mathrm{Ad}^*_g \lambda
= \lambda\,\}$ acts freely and properly on $\mu^{-1}(\lambda)$.  Then the
\emph{reduced space}
\[
  M_\lambda := \mu^{-1}(\lambda) / G_\lambda
\]
carries a unique symplectic form $\omega_\lambda$ characterised by
$\pi^* \omega_\lambda = \iota^* \omega$, where $\iota: \mu^{-1}(\lambda) \hookrightarrow M$
and $\pi: \mu^{-1}(\lambda) \twoheadrightarrow M_\lambda$.
\end{theorem}

\begin{proof}[Sketch of proof]
The tangent space to $\mu^{-1}(\lambda)$ at $x$ is $\ker d\mu_x = (\mathfrak{g}
\cdot x)^{\omega}$---the symplectic complement of the $G$-orbit through $x$.
The quotient by $G_\lambda$ removes the degenerate directions, and one verifies
that $\omega_\lambda$ is well-defined, closed, and non-degenerate.
\end{proof}

\begin{remark}
Marsden--Weinstein reduction is the symplectic analogue of quotienting by a
symmetry group.  It is used throughout mathematical physics: e.g., reducing a
mechanical system by rotation symmetry yields the effective one-dimensional
problem in the radial coordinate.
\end{remark}

\subsection{The convexity theorem}

\begin{theorem}[Atiyah 1982; Guillemin--Sternberg 1982]
If $(M, \omega)$ is compact and connected and $T$ is a torus acting on $M$
with moment map $\mu: M \to \mathfrak{t}^*$, then the image $\mu(M)$ is a
convex polytope (the \emph{moment polytope}).
\end{theorem}

% =============================================================================
\section{Gromov's Non-Squeezing Theorem}
\label{sec:gromov}
% =============================================================================

Gromov's non-squeezing theorem (1985) was a landmark discovery: it showed
that symplectic geometry is genuinely different from volume-preserving geometry.

\begin{theorem}[Gromov, 1985]\label{thm:gromov}
Let $B^{2n}(r) = \{(q,p) \in \mathbb{R}^{2n} : |q|^2 + |p|^2 < r^2\}$
be the $2n$-dimensional ball of radius $r$, and let
$Z^{2n}(R) = \{(q,p) \in \mathbb{R}^{2n} : q_1^2 + p_1^2 < R^2\}$
be the symplectic cylinder of radius $R$ (a cylinder in the first symplectic
factor $\mathbb{R}^2$).  Then a symplectic embedding
$\varphi: B^{2n}(r) \hookrightarrow Z^{2n}(R)$ exists if and only if
$r \leq R$.
\end{theorem}

\begin{remark}
By contrast, a \emph{volume-preserving} embedding $B^{2n}(r) \hookrightarrow
Z^{2n}(R)$ exists for any $R > 0$ as long as the volume fits (by a result of
Dacorogna--Moser).  Thus the non-squeezing theorem reveals a genuinely
symplectic obstruction.
\end{remark}

\begin{proof}[Idea of proof]
The key idea is the \emph{symplectic width} (or \emph{Gromov width}).  Define
the width of a symplectic manifold $(M, \omega)$ as
\[
  w(M, \omega) := \sup\bigl\{\pi r^2 : B^{2n}(r) \xhookrightarrow{\text{symp}} M\bigr\}.
\]
One establishes:
\begin{enumerate}[label=(\roman*)]
  \item The width is a symplectic invariant: $\varphi^* \omega = \omega
    \Rightarrow w(M, \omega) = w(M', \omega')$ for symplectomorphic manifolds.
  \item $w(B^{2n}(r)) = \pi r^2$.
  \item $w(Z^{2n}(R)) = \pi R^2$.
\end{enumerate}
The first item uses the theory of $J$-holomorphic curves: Gromov introduced
pseudo-holomorphic curves into symplectic topology and used the existence of
$J$-holomorphic discs in $B^{2n}(r)$ (bounded by a Lagrangian torus) to show
that any symplectic image must have a non-constant $J$-holomorphic disc
intersecting a generic $\mathbb{R}^2$-slice, whose area (by Stokes' theorem
applied to the $J$-holomorphic equation) is at least $\pi r^2$.  Since the
cylinder $Z^{2n}(R)$ has 2D-width $\pi R^2$, we obtain $r \leq R$.
\end{proof}

The non-squeezing theorem introduced the notion of \emph{symplectic capacity}:

\begin{definition}
A \emph{symplectic capacity} on the class of $2n$-dimensional symplectic
manifolds is a function $c$ assigning to each $(M, \omega)$ a value in
$[0, \infty]$ satisfying
\begin{enumerate}[label=(\roman*)]
  \item \textbf{Monotonicity:} $(M_1, \omega_1) \xhookrightarrow{\text{symp}}
    (M_2, \omega_2) \Rightarrow c(M_1) \leq c(M_2)$.
  \item \textbf{Conformality:} $c(M, \lambda\omega) = \lambda\, c(M,\omega)$.
  \item \textbf{Normalisation:} $c(B^{2n}(1)) = \pi = c(Z^{2n}(1))$.
\end{enumerate}
\end{definition}

% =============================================================================
\section{Floer Homology}
\label{sec:floer}
% =============================================================================

Floer homology, introduced by Andreas Floer in 1988--1989, is an
infinite-dimensional Morse theory on the space of loops in a symplectic
manifold.  It was created specifically to attack the Arnold conjecture.

\subsection{The action functional}

Let $(M, \omega)$ be a closed symplectic manifold and $H: S^1 \times M \to
\mathbb{R}$ a time-dependent Hamiltonian.  The \emph{symplectic action
functional} on the space $\mathcal{L}M$ of contractible loops $x: S^1 \to M$
is
\begin{equation}\label{eq:action}
  \mathcal{A}_H(x) = -\int_{D^2} \bar{x}^*\omega + \int_0^1 H(t, x(t))\,dt,
\end{equation}
where $\bar{x}: D^2 \to M$ is any disc bounding $x$.  The critical points
of $\mathcal{A}_H$ are exactly the 1-periodic orbits of $X_H$.

\subsection{Floer chain complex}

Fix a $\omega$-compatible almost complex structure $J = J(t)$.  The
\emph{Floer chain complex} is generated (over $\mathbb{Z}/2$) by the
1-periodic orbits $x$ of $X_H$, with grading given by the
\emph{Conley--Zehnder index} $\mu_{\mathrm{CZ}}(x)$.

The boundary operator $\partial: CF_k \to CF_{k-1}$ counts rigid
$J$-holomorphic cylinders (Floer trajectories) connecting orbits of adjacent
index: maps $u: \mathbb{R} \times S^1 \to M$ satisfying the
\emph{Floer equation}
\begin{equation}\label{eq:floer}
  \frac{\partial u}{\partial s} + J(t, u)\frac{\partial u}{\partial t} =
  \nabla H(t, u),
\end{equation}
with $\lim_{s \to \pm\infty} u(s,\cdot) = x^\pm(\cdot)$.

\begin{theorem}[Floer, 1989]\label{thm:floer}
$\partial^2 = 0$, so $(CF_*, \partial)$ is a chain complex.  Its homology
\[
  HF_*(M, H) := H_*(CF_*, \partial)
\]
is independent of the choices of $H$ and $J$, and is isomorphic to the
singular homology of $M$:
\[
  HF_*(M) \cong H_*(M; \mathbb{Z}/2).
\]
\end{theorem}

\subsection{Partial proof of the Arnold conjecture}

\begin{theorem}[Floer, 1989 — monotone case]\label{thm:arnold_floer}
Let $(M, \omega)$ be a monotone closed symplectic manifold
(i.e., $[\omega] = \lambda c_1(TM)$ for some $\lambda > 0$).  Then every
non-degenerate Hamiltonian diffeomorphism $\varphi \in \mathrm{Ham}(M)$
has at least $\sum_k \beta_k(M)$ fixed points.
\end{theorem}

\begin{proof}[Sketch]
The non-degenerate fixed points of $\varphi$ correspond to non-degenerate
1-periodic orbits of $H$, hence to generators of the Floer complex.  The
Euler characteristic argument gives
\[
  \#\mathrm{Fix}(\varphi) \geq \mathrm{rank}\, HF_*(M) = \mathrm{rank}\, H_*(M),
\]
which is $\sum_k \beta_k(M)$.
\end{proof}

\begin{remark}
The general Arnold conjecture (for all closed symplectic manifolds, not just
monotone ones) remains \textbf{open}.  Partial results exist for many classes
of manifolds (K\"{a}hler, Calabi--Yau, etc.), but a complete proof has not
yet been achieved.
\end{remark}

% =============================================================================
\section{Applications}
\label{sec:applications}
% =============================================================================

\subsection{Integrable systems}

A Hamiltonian system with $n$ degrees of freedom is called \emph{(Liouville)
integrable} if it admits $n$ functionally independent, pairwise Poisson-
commuting first integrals $F_1 = H, F_2, \ldots, F_n$:
\[
  \{F_i, F_j\} = 0, \quad i \neq j.
\]
The Liouville--Arnold theorem then guarantees that the common level sets
$\{F_1 = c_1, \ldots, F_n = c_n\}$ are (when compact and connected)
diffeomorphic to $n$-tori $\mathbb{T}^n$, and the motion on each torus is
quasi-periodic.

\begin{example}
The Kepler problem (gravitational two-body problem) in $\mathbb{R}^3$ is
integrable: the independent commuting integrals are $H$ (energy) and the
angular momentum components $L_1, L_2, L_3$ (with $\{L_i, L_j\} =
\varepsilon_{ijk}L_k$, but Casimir functions reduce to two commuting ones).
The Runge--Lenz vector provides an additional hidden symmetry.
\end{example}

\subsection{KAM theory}

The Kolmogorov--Arnold--Moser (KAM) theorem addresses the stability of
integrable systems under small perturbations.

\begin{theorem}[KAM; Kolmogorov 1954, Arnold 1963, Moser 1962]
Let the Hamiltonian $H_\varepsilon = H_0 + \varepsilon H_1$ be a small
perturbation of an integrable system.  Then, for sufficiently small
$\varepsilon$, a large measure of invariant tori persist---those with
\emph{Diophantine} frequency vectors $\omega = (\omega_1, \ldots, \omega_n)$
satisfying
\[
  |k \cdot \omega| \geq \frac{c}{|k|^\tau}, \quad \forall k \in \mathbb{Z}^n
  \setminus \{0\},
\]
for some $c > 0$, $\tau > n-1$.  The destroyed tori correspond to
resonances.
\end{theorem}

\subsection{Geometric quantisation}

Geometric quantisation attempts to construct a quantum theory from a
symplectic manifold $(M, \omega)$ by:
\begin{enumerate}[label=(\arabic*)]
  \item \textbf{Pre-quantisation:} Construct a Hermitian line bundle $\mathcal{L}
    \to M$ with connection $\nabla$ whose curvature is $\frac{i}{\hbar}\omega$.
    This requires the \emph{integrality condition} $[\omega/2\pi\hbar] \in
    H^2(M; \mathbb{Z})$.
  \item \textbf{Polarisation:} Choose a Lagrangian distribution (polarisation)
    to reduce the space of sections to a ``half-dimensional'' Hilbert space.
  \item \textbf{Quantisation map:} Assign to each $f \in C^\infty(M)$ a
    self-adjoint operator $\hat{f}$ satisfying $[\hat{f}, \hat{g}] = -i\hbar
    \widehat{\{f,g\}}$.
\end{enumerate}

The Hilbert space of the harmonic oscillator ($n = 1$) constructed this way
has basis $\{|0\rangle, |1\rangle, |2\rangle, \ldots\}$ with energies
$E_k = \hbar\omega(k + \frac{1}{2})$---exactly the quantum harmonic oscillator.

\subsection{Mirror symmetry}

The Homological Mirror Symmetry conjecture (Kontsevich, 1994) proposes a deep
duality between the symplectic geometry of a Calabi--Yau manifold $X$ (encoded
in its Fukaya category, whose objects are Lagrangian submanifolds) and the
complex geometry of a mirror $\hat{X}$ (encoded in its derived category of
coherent sheaves).  Floer homology groups between Lagrangian submanifolds
correspond to Ext groups of coherent sheaves.  This remains an active area
of research at the frontier of symplectic geometry and algebraic geometry.

% =============================================================================
\section*{Conclusion}
% =============================================================================

Symplectic geometry provides the mathematical framework for Hamiltonian
mechanics and has evolved into a deep theory with far-reaching implications.
The key results presented here---Darboux's theorem (local triviality),
Liouville's theorem (volume preservation), and Gromov's non-squeezing theorem
(symplectic rigidity)---together illustrate the duality between local
flexibility and global constraints that characterises the field.  Floer
homology, building on earlier work in Morse theory, provides powerful algebraic
tools for studying global properties of symplectic manifolds and partially
establishes the Arnold conjecture---one of the central open problems in the
area.

% =============================================================================
\begin{thebibliography}{99}
% =============================================================================

\bibitem{Arnold}
V.~I.~Arnold.
\newblock {\em Mathematical Methods of Classical Mechanics}.
\newblock Springer, 2nd edition, 1989.

\bibitem{CannasDS}
A.~Cannas da Silva.
\newblock {\em Lectures on Symplectic Geometry}.
\newblock Springer Lecture Notes in Mathematics 1764, 2001.

\bibitem{McDuffSalamon}
D.~McDuff and D.~Salamon.
\newblock {\em Introduction to Symplectic Topology}.
\newblock Oxford University Press, 3rd edition, 2017.

\bibitem{Gromov85}
M.~Gromov.
\newblock Pseudo holomorphic curves in symplectic manifolds.
\newblock {\em Inventiones Mathematicae}, 82:307--347, 1985.

\bibitem{Floer89}
A.~Floer.
\newblock Symplectic fixed points and holomorphic spheres.
\newblock {\em Communications in Mathematical Physics}, 120:575--611, 1989.

\bibitem{MarsdenWeinstein}
J.~Marsden and A.~Weinstein.
\newblock Reduction of symplectic manifolds with symmetry.
\newblock {\em Reports on Mathematical Physics}, 5(1):121--130, 1974.

\bibitem{ConleyZehnder}
C.~Conley and E.~Zehnder.
\newblock The Birkhoff--Lewis fixed point theorem and a conjecture of V.~I.~Arnold.
\newblock {\em Inventiones Mathematicae}, 73:33--49, 1983.

\bibitem{Darboux}
G.~Darboux.
\newblock Sur le probl\`{e}me de Pfaff.
\newblock {\em Bulletin des Sciences Math\'{e}matiques}, 6:14--36, 49--68, 1882.

\bibitem{MoserPath}
J.~Moser.
\newblock On the volume elements on a manifold.
\newblock {\em Transactions of the American Mathematical Society}, 120:286--294, 1965.

\end{thebibliography}

\end{document}
